\section{Conclusion}
This study presented an object-conditioned strategy-level parameter-bundle configuration framework for haptic teleoperation. A semantic mapping assigned each detected object class to an operational strategy that jointly specified slave-side impedance, master-side haptic rendering, and gripper execution. An asynchronous one-shot update was designed to integrate low-rate visual perception without placing inference in the nominal 200-Hz supervisory call path.

Across 27 matched task blocks comprising 135 physical trials, the complete parameter bundle reduced mean completion time by 1.79 s (8.5\%; matched-block bootstrap 95\% CI, 1.10--2.51 s) relative to impedance-only reconfiguration. Completion time favored the complete bundle in 22 of 27 matched blocks and in the mean for each operator and object, while success rate and Raw NASA-TLX were descriptively best in this mode. The post hoc sensitivity analysis restricted to matched blocks in which both trials succeeded yielded a similar mean difference of 1.80 s. No clear reduction in master-side trajectory length was observed.

These results describe a within-sample completion-time advantage of the complete parameter bundle under the tested conditions. Because the haptic-interface and gripper parameter groups were removed together in the primary comparison, their individual causal contributions remain unresolved, and the study does not demonstrate a dynamic coupling or synergy mechanism. Formal trials did not retain event-level visual histories or synchronized trigger--contact timestamps. Thus, they do not verify online trigger correctness or pre-contact completion of configuration in each trial. Future work will add synchronized trigger--contact logging, channel-specific ablations, post-contact verification, and a larger participant sample.


\section*{Supplementary Material}
Supplementary Tables S1--S3 are supplied as a separate editable supplementary file.
